% GENERATED FILE. Edit canonical lessons, then run npm run generate:books.
% canonical-source-hash: fnv1a64:893534c06ad4bd8c
% canonical-lessons: JA-C138-okimasu, JA-C138-hatarakimasu, JA-C138-arukimasu, JA-C138-kaerimasu, JA-C138-nemasu

\chapter{My Day: Get Up, Work, Walk, Go Home, Sleep}
\label{ch:ja-my-day-get-up-work-walk-go-home-sleep}
\hlchaptermodality{138}

\begin{chapteropening}
\textbf{By the end of this chapter:} \emph{I can say what I do every day with the polite -masu form, and find the stem of a verb: I get up, I work, I walk, I go home, I go to bed.}

Say a whole day in order with five polite verbs, okimasu, hatarakimasu, arukimasu, kaerimasu and nemasu, and give the dictionary form beside each one.
\end{chapteropening}

\section[\texorpdfstring{\ja{おきます} (okimasu)}{おきます (okimasu)}]{\ja{おきます} (okimasu) — I get up; to get up (polite)}

\label{lesson:JA-C138-okimasu}



\begin{quote}
Three recalls before the new word.

\begin{itemize}
\raggedright
  \item \emph{Recall:} write \textbf{\ja{ぷ}} — \textbf{R2}, five lessons back
  \item \emph{Recall:} write \textbf{\ja{ぞ}} — \textbf{R3}, twenty lessons back
  \item \emph{Recall:} say \emph{bright red} — \textbf{R4}, eighty lessons back
\end{itemize}
\end{quote}

\subsection*{You'll want to know: \ja{おきます}}
\textbf{\ja{おきます}} — \emph{okimasu} — \textquotedblleft{}I get up\textquotedblright{}, or \textquotedblleft{}she gets up\textquotedblright{}, \textquotedblleft{}we get up\textquotedblright{}: Japanese leaves out who does it when that is clear. It says what someone does, as a habit or as a plan.

Every verb in this chapter ends in \textbf{\ja{ます}}, \emph{masu}. That ending makes the polite present, the form you use with anyone you do not know well. A word list gives the dictionary form instead, here \textbf{\ja{おきる}}, \emph{okiru}. To make the polite form of a verb in \emph{-iru} or \emph{-eru} like this one, take off the last \textbf{\ja{る}} and add \textbf{\ja{ます}}: \emph{oki-ru} becomes \emph{oki-masu}.

\begin{quote}
\textbf{\ja{まいあさ} \ja{おきます。}} — \emph{maiasa okimasu} — \textquotedblleft{}I get up every morning.\textquotedblright{}
\end{quote}

\subsection*{Your turn}
\begin{itemize}
\raggedright
  \item \emph{Say it:} \emph{okimasu}
  \item \emph{Say it:} \emph{maiasa okimasu} — I get up every morning
  \item \emph{Recall:} say the dictionary form, \emph{okiru}, then the polite form, \emph{okimasu}
\end{itemize}

\begin{tcolorbox}[breakable,colback=teal!4,colframe=teal!35!black,title={Before you move on}]
How do you say \textquotedblleft{}I get up\textquotedblright{}, politely? (\textbf{\ja{おきます}}, \emph{okimasu}.) Say it once more.
\end{tcolorbox}
\section[\texorpdfstring{\ja{はたらきます} (hatarakimasu)}{はたらきます (hatarakimasu)}]{\ja{はたらきます} (hatarakimasu) — I work; to work (polite)}

\label{lesson:JA-C138-hatarakimasu}



\begin{quote}
Four recalls before the new word.

\begin{itemize}
\raggedright
  \item \emph{Recall:} say \emph{I get up}, politely — \textbf{R1}, one lesson back
  \item \emph{Recall:} say \emph{tempura} — \textbf{R2}, five lessons back
  \item \emph{Recall:} say \emph{to count} — \textbf{R3}, twenty lessons back
  \item \emph{Recall:} say \emph{vague} — \textbf{R4}, eighty lessons back
\end{itemize}
\end{quote}

\subsection*{You'll want to know: \ja{はたらきます}}
\textbf{\ja{はたらきます}} — \emph{hatarakimasu} — \textquotedblleft{}I work\textquotedblright{}: to do a job, at a desk, in a shop or in a field. The noun for the job itself is \textbf{\ja{しごと}}, which you already know.

Its dictionary form is \textbf{\ja{はたらく}}, \emph{hataraku}. Verbs that end in \emph{-u} like this one change that last sound to \emph{-i} before \textbf{\ja{ます}}: \emph{hatarak-u} becomes \emph{hatarak-i-masu}. So \textbf{\ja{く}} turns into \textbf{\ja{き}}.

\begin{quote}
\textbf{\ja{ぎんこうで} \ja{はたらきます。}} — \emph{ginkō de hatarakimasu} — \textquotedblleft{}I work at a bank.\textquotedblright{} The particle \textbf{\ja{で}} marks the place where something is done.
\end{quote}

\subsection*{Your turn}
\begin{itemize}
\raggedright
  \item \emph{Say it:} \emph{hatarakimasu}
  \item \emph{Say it:} \emph{maiasa okimasu}, then \emph{ginkō de hatarakimasu}
  \item \emph{Recall:} change \emph{hataraku} to its polite form, and say it
\end{itemize}

\begin{tcolorbox}[breakable,colback=teal!4,colframe=teal!35!black,title={Before you move on}]
How do you say \textquotedblleft{}I work\textquotedblright{}, politely? (\textbf{\ja{はたらきます}}, \emph{hatarakimasu}.) Say it once more.
\end{tcolorbox}
\section[\texorpdfstring{\ja{あるきます} (arukimasu)}{あるきます (arukimasu)}]{\ja{あるきます} (arukimasu) — I walk; to walk (polite)}

\label{lesson:JA-C138-arukimasu}



\begin{quote}
Four recalls before the new word.

\begin{itemize}
\raggedright
  \item \emph{Recall:} say \emph{I work}, politely — \textbf{R1}, one lesson back
  \item \emph{Recall:} say \emph{fluent} — \textbf{R2}, five lessons back
  \item \emph{Recall:} say \emph{water} — \textbf{R3}, twenty lessons back
  \item \emph{Recall:} say \emph{certain, sure} — \textbf{R4}, eighty lessons back
\end{itemize}
\end{quote}

\subsection*{You'll want to know: \ja{あるきます}}
\textbf{\ja{あるきます}} — \emph{arukimasu} — \textquotedblleft{}I walk\textquotedblright{}, \textquotedblleft{}I go on foot\textquotedblright{}.

The dictionary form is \textbf{\ja{あるく}}, \emph{aruku}, and it works like \textbf{\ja{はたらく}}: the last \textbf{\ja{く}} becomes \textbf{\ja{き}}, and \textbf{\ja{ます}} follows. \emph{aruk-u} becomes \emph{aruk-i-masu}.

\begin{quote}
\textbf{\ja{えきまで} \ja{あるきます。}} — \emph{eki made arukimasu} — \textquotedblleft{}I walk to the station.\textquotedblright{} \textbf{\ja{まで}} means \textquotedblleft{}as far as\textquotedblright{}.
\end{quote}

\subsection*{Your turn}
\begin{itemize}
\raggedright
  \item \emph{Say it:} \emph{arukimasu}
  \item \emph{Say it:} \emph{eki made arukimasu} — I walk to the station
  \item \emph{Recall:} say \emph{I work}, then \emph{I walk}, and listen for the \emph{-kimasu} in both
\end{itemize}

\begin{tcolorbox}[breakable,colback=teal!4,colframe=teal!35!black,title={Before you move on}]
How do you say \textquotedblleft{}I walk\textquotedblright{}, politely? (\textbf{\ja{あるきます}}, \emph{arukimasu}.) Say it once more.
\end{tcolorbox}
\section[\texorpdfstring{\ja{かえります} (kaerimasu)}{かえります (kaerimasu)}]{\ja{かえります} (kaerimasu) — I go home, I go back; to return (polite)}

\label{lesson:JA-C138-kaerimasu}



\begin{quote}
Three recalls before the new word.

\begin{itemize}
\raggedright
  \item \emph{Recall:} say \emph{I walk}, politely — \textbf{R1}, one lesson back
  \item \emph{Recall:} write \textbf{\ja{ぺ}} — \textbf{R2}, five lessons back
  \item \emph{Recall:} write \textbf{\ja{ず}} — \textbf{R3}, twenty lessons back
\end{itemize}
\end{quote}

\subsection*{You'll want to know: \ja{かえります}}
\textbf{\ja{かえります}} — \emph{kaerimasu} — \textquotedblleft{}I go home\textquotedblright{}, \textquotedblleft{}I go back\textquotedblright{}, to the place you belong to: your home, your town, your country.

Its dictionary form is \textbf{\ja{かえる}}, \emph{kaeru}. It ends in \emph{-eru}, but it follows the \emph{-u} pattern of \textbf{\ja{あるく}}, not the pattern of \textbf{\ja{おきる}}: the last \textbf{\ja{る}} becomes \textbf{\ja{り}}, and \textbf{\ja{ます}} follows. So it is \emph{kaer-i-masu}, never \emph{kaemasu}. A few common verbs in \emph{-eru} and \emph{-iru} work this way, and a word list marks them.

\begin{quote}
\textbf{\ja{うちに} \ja{かえります。}} — \emph{uchi ni kaerimasu} — \textquotedblleft{}I go home.\textquotedblright{} The particle \textbf{\ja{に}} marks where you are going.
\end{quote}

\subsection*{Your turn}
\begin{itemize}
\raggedright
  \item \emph{Say it:} \emph{kaerimasu}
  \item \emph{Say it:} \emph{uchi ni kaerimasu} — I go home
  \item \emph{Recall:} say \emph{I walk}, then \emph{I go home}
\end{itemize}

\begin{tcolorbox}[breakable,colback=teal!4,colframe=teal!35!black,title={Before you move on}]
How do you say \textquotedblleft{}I go home\textquotedblright{}, politely? (\textbf{\ja{かえります}}, \emph{kaerimasu}.) Say it once more.
\end{tcolorbox}
\section[\texorpdfstring{\ja{ねます} (nemasu)}{ねます (nemasu)}]{\ja{ねます} (nemasu) — I go to bed, I sleep; to sleep (polite)}

\label{lesson:JA-C138-nemasu}



\begin{quote}
Two recalls before the new word.

\begin{itemize}
\raggedright
  \item \emph{Recall:} say \emph{I go home}, politely — \textbf{R1}, one lesson back
  \item \emph{Recall:} say \emph{cool} — \textbf{R3}, twenty lessons back
\end{itemize}
\end{quote}

\subsection*{You'll want to know: \ja{ねます}}
\textbf{\ja{ねます}} — \emph{nemasu} — \textquotedblleft{}I go to bed\textquotedblright{}, \textquotedblleft{}I sleep\textquotedblright{}.

Its dictionary form is \textbf{\ja{ねる}}, \emph{neru}, and it works like \textbf{\ja{おきる}}: take off \textbf{\ja{る}} and add \textbf{\ja{ます}}. \emph{ne-ru} becomes \emph{ne-masu}.

\begin{quote}
\textbf{\ja{いま} \ja{ねます。}} — \emph{ima nemasu} — \textquotedblleft{}I am going to bed now.\textquotedblright{}
\end{quote}

Now you have a whole day, and the two ways to find the stem:

\noindent
\begin{tabularx}{\linewidth}{@{}>{\raggedright\arraybackslash}X>{\raggedright\arraybackslash}X>{\raggedright\arraybackslash}X@{}}
\toprule
\textbf{dictionary form} & \textbf{polite present} & \textbf{how} \\
\midrule
\textbf{\ja{おきる}} \emph{okiru} & \textbf{\ja{おきます}} \emph{okimasu} & drop \emph{-ru} \\
\textbf{\ja{ねる}} \emph{neru} & \textbf{\ja{ねます}} \emph{nemasu} & drop \emph{-ru} \\
\textbf{\ja{はたらく}} \emph{hataraku} & \textbf{\ja{はたらきます}} \emph{hatarakimasu} & \emph{-u} becomes \emph{-i} \\
\textbf{\ja{あるく}} \emph{aruku} & \textbf{\ja{あるきます}} \emph{arukimasu} & \emph{-u} becomes \emph{-i} \\
\textbf{\ja{かえる}} \emph{kaeru} & \textbf{\ja{かえります}} \emph{kaerimasu} & \emph{-u} becomes \emph{-i} \\
\bottomrule
\end{tabularx}

\subsection*{Your turn}
\begin{itemize}
\raggedright
  \item \emph{Say it:} \emph{nemasu}
  \item \emph{Say it:} your day in order, politely: \emph{okimasu}, \emph{hatarakimasu}, \emph{arukimasu}, \emph{kaerimasu}, \emph{nemasu}
  \item \emph{Recall:} for each of the five, say the dictionary form, then the polite form
\end{itemize}

\begin{tcolorbox}[breakable,colback=teal!4,colframe=teal!35!black,title={Before you move on}]
I get up? (\textbf{\ja{おきます}}, \emph{okimasu}.) I work? (\textbf{\ja{はたらきます}}, \emph{hatarakimasu}.) I walk? (\textbf{\ja{あるきます}}, \emph{arukimasu}.) I go home? (\textbf{\ja{かえります}}, \emph{kaerimasu}.) I go to bed? (\textbf{\ja{ねます}}, \emph{nemasu}.)
\end{tcolorbox}
